\section{Perturbative and Tensorial Framework}

In this section we introduce the perturbative and tensorial framework
used in the present work at the level required for the main text.
The aim is not to rederive rovibrational perturbation theory in full
detail, but to clarify the structure of the tensorial objects entering
the centrifugal distortion problem and their relation to the
quantities obtained in standard VPT-based calculations. Detailed
derivations of the quartic and sextic tensor constructions, explicit
projection formulas, and full channel algebra are deferred to the
appendices.

The key point is that the present framework is built around two
logically distinct layers. The first is the standard linear sector, in
which quartic and standard sextic centrifugal distortion can be lifted
from Watson constants to tensor representatives, transported between
representations, and reprojected onto the desired Hamiltonian
parametrization. The second is the first breakdown of that picture,
which occurs when new tensorial objects appear that cannot be absorbed
into the first-level linear sector.

At the most basic conceptual level, the framework may be read as
follows. Watson constants are not treated as primitive quantities, but
as projected coordinates of underlying tensor representatives. These
representatives are reconstructed by pseudoinverse gauge fixing,
transported linearly between representations within the standard
sector, and then reprojected onto the desired Hamiltonian
parametrization. The main theoretical problem is therefore to identify
which contributions belong to this minimal linear sector and which ones
do not.

This distinction should not be read as a mere change of coordinates in
the parameter space of the Watson Hamiltonian. The purpose of the
tensorial formulation is to identify the minimal linear transport
sector underlying the standard quartic and sextic problems and then to
locate the first contribution that lies outside it.

In operational terms, the standard problem is organized by a simple
sequence:
\begin{equation}
\text{Watson constants}
\;\longrightarrow\;
\text{tensor representative}
\;\longrightarrow\;
\text{transport}
\;\longrightarrow\;
\text{reprojection}.
\end{equation}
Conceptually, the procedure therefore involves three steps:
reconstruction of tensor representatives from Watson constants, linear
transport of these tensors under representation or axis changes, and
projection onto the desired Hamiltonian parametrization.
Because the inverse step is rank-deficient, the tensor representative
is selected by Moore--Penrose pseudoinversion and should be understood
as a gauge-fixed representative rather than as a uniquely defined
object, with the pseudoinverse providing a natural gauge fixing of the
reconstruction problem.\cite{Penrose1955} The standard quartic and sextic sectors are precisely those for
which this gauge-fixed transport remains linear.

The purpose of the present section is therefore threefold: to specify
the perturbative partition underlying the channel language, to identify
the tensorial objects spanning the standard linear sector, and to show
how the first quartic and sextic diagnostics of breakdown arise once
that sector ceases to be sufficient.
The derivation therefore proceeds by identifying the minimal tensor
objects required to generate centrifugal distortion constants and by
analyzing how their projections define the parameters of the effective
Hamiltonian.

For readers approaching the problem from standard rotational
spectroscopy rather than from tensor methods, the logical progression
of the section is straightforward. We first define the perturbative
partition, then identify the tensorial objects generating the standard
quartic and standard sextic sectors, and finally show how the first
genuinely new tensor structures appear once one leaves this minimal
linear sector.

\begin{table}[h]
\centering
\caption{Minimal notation used in the main text.}
\begin{tabular}{ll}
\hline
Symbol & Meaning \\
\hline
Watson constants & Quartic or sextic effective-Hamiltonian parameters \\
Tensor representative & Gauge-fixed tensor object reconstructed from Watson constants \\
Transported tensor & Tensor representative after axis/reduction transport \\
$\mu_1$ & First derivative of the inverse inertia tensor with respect to $Q$ \\
$\tau$ & Standard quartic tensor entering quartic/sextic transport \\
$H_{22}$ & First quartic channel outside the minimal linear sector \\
\hline
\end{tabular}
\end{table}

\subsection{Perturbative framework and partition convention}

The rovibrational Hamiltonian is expanded in terms of vibrational
coordinates and rotational operators. Following standard notation,
individual contributions may be grouped into blocks $H_{nm}$, where
$n$ denotes the vibrational order and $m$ the rotational order.
These labels identify the vibrational--rotational bidegree of the
individual terms but do not, by themselves, determine the perturbative
order. The latter is defined only after the Van Vleck contact
transformation is applied and the Baker--Campbell--Hausdorff expansion
is evaluated, which generates the familiar combinatorial prefactors
and factorial denominators of rovibrational perturbation
theory.\cite{Watson1967,Watson1968,Watson1968JCP,AlievWatson1976}

In the present work the standard Watson/VPT2-style partition is used.
Within this convention,
\begin{equation}
H_{03},\; H_{12},\; H_{30}
\end{equation}
are treated as first-order terms, whereas
\begin{equation}
H_{40},\; H_{22},\; H_{13},\; H_{31},\; H_{04}
\end{equation}
are treated as second-order terms. This convention is standard in the
present context, but it is not universal. In alternative formulations,
including state-summation schemes that are also often denoted as VPT2,
the same blocks may be assigned to a different effective perturbative
order.\cite{Franke2023}

Within a fixed perturbative partition and truncation order, the
effective quartic and sextic rotational Hamiltonians may be decomposed
into channels such as
\begin{equation}
H_{12}H_{12},\qquad H_{22},\qquad H_{12}H_{30},\qquad
H_{30}H_{30},\qquad H_{04},
\end{equation}
and analogous higher-order sextic contributions. This channel
decomposition is exact only within the adopted perturbative framework,
that is, after fixing the partition, the truncation order, and the
effective-Hamiltonian construction. The symbolic channel labels do not
by themselves determine the numerical weights of the corresponding
terms. Once the Van Vleck contact transformation is expanded, the
Baker--Campbell--Hausdorff series generates the usual combinatorial
prefactors and factorial denominators. Accordingly, the exact
weighting of each channel is the one induced by the chosen
perturbative construction, not by the labels $H_{nm}$ taken in
isolation.

Table~\ref{tab:channel_hierarchy_summary} summarizes the practical
organization adopted in the present work. The important distinction is
not simply between different perturbative channels, but between those
contributions that remain inside the minimal linear tensor sector and
those that introduce genuinely new tensorial structures. In this
sense, the table should be read as the Watson/VPT2 language recast in
tensor-hierarchy form.\cite{Watson1968,Watson1977,Mills1972}

\begin{table*}[t]
\centering
\caption{Summary of the main quartic and sextic contributions discussed in the
present work, organized by perturbative source, primitive tensorial content,
and role with respect to the minimal linear tensor sector.}
\begin{tabular}{lllll}
\toprule
Sector & Perturbative source / channel & Primitive tensor objects & In linear sector? & Main role \\
\midrule
Standard quartic & $H_{12}H_{12}$ & $\mu_1$, Coriolis & Yes & Standard quartic distortion constants \\
Standard sextic geometry & Sextic geometry closure in $c_1$, $c_2$, $\tau$ & $\mu_1$, Coriolis, $\tau$ & Yes & Standard sextic geometry sector \\
Standard sextic anharmonicity & Cubic-force contribution & $\phi_3$ & First nonlinear layer & Standard sextic anharmonic contribution \\
First quartic breakdown & $H_{22}$ & $\mu_2 = B(\mu_1)-N$ & No & First breakdown diagnostic for quartics \\
First sextic breakdown & Linear response in $\tau(H_{22})$ & $\tau(H_{22})$ & No & First breakdown diagnostic for sextics \\
\bottomrule
\end{tabular}
\label{tab:channel_hierarchy_summary}
\end{table*}

\subsection{Physical meaning of the tensor representatives}

Within this perturbative framework the quantities entering centrifugal
distortion theory can be organized in terms of tensorial objects
constructed from derivatives of the inertia tensor and from the usual
rovibrational coupling terms. In particular, the first derivative of
the inverse inertia tensor with respect to vibrational coordinates,
\begin{equation}
\mu_1 = \frac{\partial I^{-1}}{\partial Q},
\end{equation}
plays a central role in the construction of the standard quartic
centrifugal distortion sector.
The tensor $\mu_1$ represents the first derivative of the inverse
inertia tensor with respect to vibrational coordinates and therefore
describes the leading geometric coupling between rotational motion and
vibrational displacements. In the present formulation it plays a
central role, since it generates the minimal tensorial structure
underlying quartic centrifugal distortion constants and forms the basic
building block of the linear tensor sector discussed below. Within
this framework, higher-order distortion constants arise from
combinations and projections of these first-level tensor objects,
whereas genuinely new tensor structures appear only when anharmonic
force constants enter the expansion.

The tensor representatives used in the present work are therefore
directly connected to the quantities appearing in standard
rovibrational perturbation theory. They should not be viewed as
auxiliary algebraic encodings layered on top of VPT2, but rather as
a compact tensorial representation of the same physical objects that
enter conventional derivations of centrifugal distortion constants.
They are therefore not directly measurable quantities in the same sense
as transition frequencies, but gauge-fixed representatives of the
underlying rovibrational content encoded in the effective Hamiltonian.
In the harmonic quartic problem, the basic object is precisely the
first derivative of the inverse inertia tensor or, equivalently in the
standard harmonic setting, its Coriolis-coupling description.

More explicitly, if
\begin{equation}
\mu = I^{-1},
\end{equation}
then differentiation of the identity $I^{-1}I=\mathbf 1$ with respect to
a normal coordinate $Q_k$ gives
\begin{equation}
\frac{\partial I^{-1}}{\partial Q_k}
=
-\,I^{-1}\frac{\partial I}{\partial Q_k}I^{-1}.
\end{equation}
The standard quartic tensor is then generated by the familiar
second-order perturbative contraction
\begin{equation}
\tau_{\alpha\beta\gamma\delta}
=
-\frac12\sum_k
\frac{
\mu_{\alpha\beta}^{(k)}\mu_{\gamma\delta}^{(k)}
}{\omega_k^2},
\end{equation}
where $\mu_{\alpha\beta}^{(k)}=\partial\mu_{\alpha\beta}/\partial Q_k$.

The same quartic content may be written in Coriolis form. If
\begin{equation}
G_{k\mu}=\sum_i m_i\left(\mathbf r_i\times \mathbf l_{ik}\right)_\mu
\end{equation}
denotes the vibrational angular-momentum vector associated with mode
$k$, the corresponding symmetric Coriolis tensor is
\begin{equation}
M_{\mu\nu}=\sum_k\frac{G_{k\mu}G_{k\nu}}{\omega_k^2}.
\end{equation}
For the standard quartic contribution, the reduced tensor
representative $\tau'$ is linearly related to this Coriolis tensor.
Thus the passage from inertia derivatives to Coriolis couplings is not
an additional approximation, but an alternative tensorial
representation of the same standard quartic rovibrational content.

More concretely, the standard quartic tensor is generated by the
harmonic dependence of the inverse inertia tensor on the normal
coordinates. In the usual force-field language this means that the
quartic distortion constants are determined by the equilibrium
geometry, the Cartesian Hessian, and the associated harmonic normal
modes. In the present implementation these are the same data from
which the first inertia derivatives, Coriolis couplings, and harmonic
quartic tensor are constructed.

The same physical interpretation extends to the standard sextic
sector. The standard sextic formulas involve the quartic tensor,
Coriolis couplings, and the explicit anharmonic cubic force
constants. The important point established here is that the standard
sextic geometry sector does not require second derivatives of the
inertia tensor as primitive objects. Instead, it closes on the same
first-level geometric language that already governs the standard
quartic sector, with the cubic force constants entering only as the
explicit anharmonic contribution.

In the formulation adopted here, this statement is made explicit by
writing the sextic standard sector in terms of the working quantities
$c_1$, $c_2$, and $\tau$, which organize the geometric part of the
problem, plus a separate cubic-force contribution. At the structural
level one may therefore write
\begin{equation}
\Phi^{(6)} = \Phi^{(6)}_{\mathrm{geom}}(c_1,c_2,\tau)
+ \Phi^{(6)}_{\mathrm{cubic}}(\phi_3).
\end{equation}
This is already more explicit than the usual black-box presentation of
the sextic formulas, because it isolates the tensorial geometry from
the explicit anharmonic input.

For the present purposes it is also useful to refine the cubic term as
\begin{equation}
\Phi^{(6)}_{\mathrm{cubic}}
=
\Phi^{(6)}_{\mathrm{cubic,sd}}
+
\Phi^{(6)}_{\mathrm{cubic,3ind}},
\end{equation}
where the first piece collects the semi-diagonal cubic-force
contributions and the second the genuinely three-index remainder. This
separation is interesting in its own right, because it aligns the
sextic theory with the computational hierarchy that is realistic for
medium-sized molecules.

Schematically, the standard sextic structure is organized by
\begin{equation}
\mu_1,\qquad \zeta,\qquad \phi_3,
\end{equation}
where $\zeta$ denotes Coriolis coupling data and $\phi_3$ the cubic
force constants. In this sense, standard quartics and standard sextics
belong to the same physical layer of rovibrational perturbation
theory: they differ in the way the same first-level geometric objects
are combined, not in the primitive tensorial level required to
describe them.
The present formulation therefore reveals a natural hierarchy: harmonic
geometry defines the inertia tensor, the first derivative $\mu_1$
generates the linear tensor sector responsible for standard quartic and
standard sextic distortion, and cubic force constants introduce the
first genuinely nonlinear layer.

Accordingly, the present standard-sector construction should not be
read as an alternative to Watson/VPT2 perturbation theory, but as a
tensorial reformulation of its usual content. The required inputs are
the same ones used in standard force-field based treatments:
equilibrium geometry, Cartesian Hessian, normal coordinates, Coriolis
data, and, for the sextic standard sector, the cubic force constants
entering the explicit anharmonic contribution.

The same structural content may also be written in coordinate
languages different from the Cartesian-normal-mode one used here. In
particular, recent treatments based on internal-coordinate
formulations of vibrational and rovibrational corrections employ a
different working notation, but at the harmonic and standard-VPT2
levels the underlying rovibrational structure is the same.\cite{cubic-internal_2024}
From the present point of view, this means that the tensorial
hierarchy is not tied to a specific coordinate choice: what changes is
only the parametrization of the intermediate quantities, not the
first-level geometric content. In particular, the same families of
quantities that enter vibrational corrections to the rotational
constants (the usual $\alpha$ terms) also provide the natural building
blocks of the standard sextic structure once written in the present
tensor language.

In the curvilinear-internal-coordinate formulation this equivalence
appears already at the level of the generalized coordinate set,
\begin{equation}
\xi = \{s_1,\ldots,s_N,\phi_1,\ldots,\phi_{N_r},
X_{\mathrm{COM}},Y_{\mathrm{COM}},Z_{\mathrm{COM}}\},
\end{equation}
where the $s_i$ are internal vibrational coordinates and the $\phi_i$
describe the body-fixed orientation.\cite{cubic-internal_2024} In
that language, second-order vibrational and rovibrational corrections
are built from the same core ingredients that appear here in the
Cartesian-normal-mode formulation, namely harmonic frequencies,
Coriolis couplings, and semidiagonal third derivatives of the force
field. The harmonic rotation-vibration block is therefore the same
object in the two formulations, merely written in different
coordinates. The present tensor language should thus be read as a
coordinate-independent structural interpretation of the same
first-level rovibrational information, rather than as a separate
formalism alongside the internal-coordinate treatment.

More explicitly, in the internal-coordinate formulation the
rovibrational Hamiltonian is written in block form as
\begin{equation}
H_{\mathrm{vr}}
=
\frac12 \sum_{\tau\tau'} J_\tau g^{r}_{\tau\tau'} J_{\tau'}
+ \frac12 \sum_{i\tau}\left(J_\tau g^{vr}_{\tau i} p_i + p_i g^{vr}_{i\tau} J_\tau\right)
+ \frac12 \sum_{ij} p_i g^{v}_{ij} p_j
+ g + V,
\end{equation}
where $g^r$, $g^{vr}$, and $g^v$ are the rotational,
rotation-vibration, and vibrational blocks of the rovibrational
metric.\cite{cubic-internal_2024} This is the same harmonic
content that, in the Cartesian-normal-mode formulation used in the
present paper, appears through the inverse inertia tensor, its first
derivative $\mu_1$, and the corresponding Coriolis couplings.
The equivalence is therefore structural: the rotational block
corresponds to the inertia tensor itself, the mixed block to the same
first-order rotation-vibration coupling encoded here by $\mu_1$ and
$\zeta$, and the vibrational block to the harmonic metric defining the
normal modes. The tensor formulation adopted in the present work is
thus best understood as a coordinate-independent rewriting of the same
harmonic rovibrational block rather than as an alternative formal
construction.

Within this scope, the present formulation is intended to be
equivalent to the standard perturbative theory rather than an
additional approximation to it. The tensor representatives reorganize
the same standard content into a transport-friendly language; they do
not discard standard VPT2 terms from the quartic or standard sextic
sector.

\subsection{Pseudoinverse reconstruction as gauge fixing}

Once one adopts the tensor representatives as primary objects, the
inverse problem becomes conceptually transparent. In the quartic case,
the reduced tensor representative contains six independent
coordinates, whereas Watson's quartic Hamiltonian contains only five
spectroscopic constants. The map from tensor space to spectroscopic
space is therefore rank-deficient and the inverse reconstruction of
the tensor from a given parameter set is not unique.

This is not a pathology of the method. It is a structural fact: the
spectroscopic constants determine an affine family of tensor
representatives rather than a unique tensor. The Moore--Penrose
pseudoinverse is used here to select one distinguished representative
within that family. In other words, the pseudoinverse is a
gauge-fixing prescription for the inverse problem. More specifically,
it selects the canonical minimum-norm representative of the
gauge-equivalence class defined by the fitted constants.

Equivalently, Watson constants do not label a point in tensor space,
but a gauge orbit inside it. The tensor representative used in the
present work is therefore not ``the'' tensor behind a fitted parameter
set in an absolute sense, but the gauge-fixed representative selected
by the Moore--Penrose slice. Linear transport between representations
is to be understood precisely on this gauge-fixed standard sector.

This point matters both conceptually and numerically. Conceptually, it
separates the intrinsic tensorial content of centrifugal distortion
from the particular spectroscopic coordinates used to describe it.
Numerically, it makes explicit that the reconstruction step may become
delicate when the spectroscopic parameters are strongly correlated or
when the transformation matrix is poorly conditioned. This is
particularly relevant in real spectroscopic fits, where the fitted
constants are not statistically independent and representation changes
may amplify parameter correlations.

For the quartic problem, the relevant projection has rank five on a
six-dimensional reduced tensor space, which is precisely why the
inverse problem defines an affine gauge family rather than a unique
tensor. In the sextic problem, the transport is carried out on the
canonical physical subspace selected by the model, with the
complementary residual monitored explicitly. The numerical stability
of the reconstruction is therefore not assumed, but checked through
condition numbers, residual norms, and related diagnostic quantities.

For this reason, the present framework always couples pseudoinverse
reconstruction to explicit conditioning diagnostics. The same tensor
transport that implements the representation change also exposes
condition numbers, spectral invariants, residual norms, and related
quantities such as $s_{111}$. The reconstruction should therefore be
read not as a black-box inversion, but as a controlled numerical gauge
choice whose stability can be monitored.

Once the tensor representative has been reconstructed, representation
and reduction changes are carried out by:
\begin{enumerate}
\item lifting the Watson constants to the tensor representative;
\item transporting the tensor under the corresponding axis permutation;
\item reprojecting the transported tensor onto the target effective-Hamiltonian
coordinates.
\end{enumerate}

This lift--transport--reproject structure is the basic mechanism used
throughout the standard quartic and sextic sectors. The operational
logic of the framework is summarized in Fig.~\ref{fig:ceditt3_scheme}.
Standard quartic and sextic constants are lifted to tensor
representatives by pseudoinverse reconstruction, transported under
representation changes, and finally reprojected onto the target
Watson parametrization. When harmonic input is available, the same
workflow also provides the first post-standard diagnostics.

\begin{figure}[t]
\centering
\resizebox{\textwidth}{!}{
\begin{tikzpicture}[
box/.style={
rectangle,draw,rounded corners,
align=center,
minimum width=4.3cm,
minimum height=1.3cm
},
arrow/.style={->,thick}
]

% INPUTS
\node[box] (fit)
{Spectroscopic fit\\Watson constants\\$(D,H)$};

\node[box,below=1.6cm of fit] (harm)
{Harmonic input\\geometry + Hessian\\or `.fchk`};

\node[above=0.6cm of fit] {\large \textbf{Input}};

% LIFT
\node[box,right=4.9cm of fit] (lift)
{Lift / reconstruction\\pseudoinverse gauge fixing\\tensor representatives};

\node[above=0.6cm of lift] {\large \textbf{Lift}};

% TRANSPORT
\node[box,right=4.9cm of lift] (transport)
{Transport\\axis permutation\\reduction / representation change};

\node[above=0.6cm of transport] {\large \textbf{Transport}};

% OUTPUT
\node[box,right=4.9cm of transport,yshift=1.35cm] (reproj)
{Standard output\\reprojected Watson constants\\in chosen parametrization};

\node[box,right=4.9cm of transport,yshift=-1.35cm] (diag)
{Post-standard diagnostics\\conditioning, $s_{111}$, $H_{22}$,\\linear sextic $\tau(H_{22})$ term};

\node[above=0.6cm of reproj] {\large \textbf{Output}};

% ARROWS
\draw[arrow] (fit.east) -- (lift.west);
\draw[arrow] (harm.east) -- (diag.west);
\draw[arrow] (lift.east) -- (transport.west);
\draw[arrow] (transport.east) -- (reproj.west);
\draw[arrow] (transport.east) -- (diag.west);

\end{tikzpicture}
}

\caption{
Operational structure of the present framework. Standard quartic and sextic
centrifugal distortion constants are lifted from Watson parametrizations to
tensor representatives through pseudoinverse gauge fixing, transported under
representation changes, and reprojected onto the target Hamiltonian
parametrization. Optional harmonic input is only needed for post-standard
diagnostics, namely the quartic $H_{22}$ contribution and its first sextic
linear analogue.
}

\label{fig:ceditt3_scheme}
\end{figure}


The conceptual hierarchy behind this workflow is summarized separately
in Fig.~\ref{fig:ceditt3_hierarchy}. Harmonic geometry first defines
the inertia tensor, the first derivative $\mu_1$ generates the minimal
linear tensor sector shared by standard quartic and standard sextic
centrifugal distortion, and cubic force constants introduce the first
genuinely nonlinear layer.

\begin{figure}[t]
\centering
\resizebox{0.78\textwidth}{!}{
\begin{tikzpicture}[
box/.style={
rectangle,draw,rounded corners,
align=center,
minimum width=5.4cm,
minimum height=1.15cm
},
arrow/.style={->,thick}
]

\node[box] (harm) {Harmonic geometry\\equilibrium structure + inertia tensor};
\node[box,below=1.3cm of harm] (muone) {$\mu_1=\partial I^{-1}/\partial Q$\\first geometric rotation--vibration coupling};
\node[box,below=1.3cm of muone] (linear) {Linear tensor sector\\standard quartic + standard sextic};
\node[box,below=1.3cm of linear] (cubic) {Cubic force constants\\explicit anharmonic input};
\node[box,below=1.3cm of cubic] (nonlinear) {First nonlinear layer\\first breakdown of linear transport};

\draw[arrow] (harm.south) -- (muone.north);
\draw[arrow] (muone.south) -- (linear.north);
\draw[arrow] (linear.south) -- (cubic.north);
\draw[arrow] (cubic.south) -- (nonlinear.north);

\end{tikzpicture}
}

\caption{
Conceptual hierarchy revealed by the tensor formulation. Harmonic
geometry defines the inertia tensor, its first derivative $\mu_1$
generates the minimal linear tensor sector governing standard quartic
and standard sextic centrifugal distortion constants, and cubic force
constants introduce the first genuinely nonlinear layer.
}

\label{fig:ceditt3_hierarchy}
\end{figure}


\subsection{Standard quartic sector}

In the standard quartic sector, the relevant rovibrational content is
generated entirely by the first-level inverse-inertia derivative
$\mu_1$, or, equivalently, by its Coriolis-coupling description. The
standard quartic contribution therefore remains fully inside the
first-level tensorial language.

At the spectroscopic level, the quartic constants are obtained after
projection of the quartic tensor representative onto Watson's operator
basis. At the tensor level, however, the standard quartic problem is
linear: once the representative has been fixed by pseudoinverse
reconstruction, representation changes are implemented by linear
transport of the tensor itself. This is the first half of the
standard linear sector developed in the present work.

\subsection{Standard sextic sector}

The sextic problem is structurally richer, but in its standard form it
still closes on the same first-level tensorial language. The key
result of the present analysis is that the standard sextic geometry
sector does not require second derivatives of the inertia tensor as
primitive objects. Instead, the standard sextic geometry contributions
can be written in terms of the same first-level data that govern the
standard quartic problem, together with Coriolis couplings and the
explicit anharmonic cubic force constants.

This is physically important. If second inertia derivatives were
genuinely primitive already at the level of standard sextics, then
quartics and sextics would not belong to a common linear transport
sector. By showing that the standard sextic geometry sector can be
rewritten without primitive $d^2I/dQ^2$, the present framework
identifies the correct baseline from which genuinely new
post-standard tensorial objects can later be recognized. This is the
second half of the standard linear sector.

In this sense, the ``standard sextic sector'' should be understood
strictly as the sextic content that can still be expressed in terms of
the same first-level tensorial objects governing the quartic problem.
The explicit cubic-force contribution is then precisely the first place
where this minimal linear closure must be enlarged.

\subsection{Unified standard linear sector}

The standard quartic and standard sextic problems therefore admit a
common tensorial description. They are both governed by the same
first-level geometric language, and both are transported between
representations by the same pseudoinverse-lift, tensor-transport, and
reprojection mechanism.

This observation provides a unified tensor interpretation of the
standard quartic and sextic sectors and explains why their
representation dependence can be treated within the same transport
framework.

This motivates the central notion of the present paper: there is a
standard linear representation-equivalent sector of centrifugal
distortion theory. In this sector, fitted Watson constants and
directly computed rovibrational quantities can be compared on the same
tensorial footing, without introducing additional primitive objects
beyond those already present in the standard construction.

The corresponding equivalence should be understood precisely at the
level of this minimal linear sector. It is this restricted but
physically important domain, rather than the full higher-order
effective Hamiltonian, that is shown here to admit a common tensorial
transport picture.

In practical quantum-chemical calculations the standard perturbative
sector typically includes only the $H_{12}H_{12}$ contributions. Under
these conditions the linear transport framework remains valid, and the
breakdown discussed here should be interpreted mainly as a formal
limitation of the representation-equivalence structure rather than as a
practical computational difficulty.

\subsection{First breakdown in the quartic sector: \texorpdfstring{$H_{22}$}{H22}}

The first departure from this standard linear sector occurs in the
quartic channel $H_{22}$. This channel is structurally simple enough
to be isolated, yet it already forces the appearance of a genuinely
new tensorial object.

In the present work, the expression ``first breakdown'' has a precise
meaning. It denotes the first contribution to the effective
centrifugal-distortion problem that cannot be represented within the
minimal first-level linear tensor sector generated by $\mu_1$ and its
standard geometric descendants. Operationally, it is therefore the
first term that makes linearly transported Watson parameter sets in
different representations cease to be strictly equivalent descriptions
of the same fitted content.

Physically, $H_{22}$ is the first contribution that depends on the
second response of the inertia tensor to vibrational motion. In the
notation adopted here, the full second inverse-inertia derivative may
be written as
\begin{equation}
\mu_2 = B(\mu_1) - N,
\end{equation}
where $B(\mu_1)$ is the bilinear closure generated algebraically from
the first-level object $\mu_1$, while $N$ is the intrinsic second-level
tensor not reducible to the standard first-level sector.

This decomposition is the physical core of the breakdown. The term
$B(\mu_1)$ does not introduce a new tensorial layer; it is the
algebraic closure of the already known first-level geometry. By
contrast, $N$ represents the genuinely new part of the second-order
inertia response. The quartic channel $H_{22}$ is therefore the first
place where the standard linear language is no longer sufficient by
itself.

In this sense, $H_{22}$ is not merely an additional post-standard
correction. It is the first identifiable tensorial marker of a new
coupling mechanism: the first point at which the effective quartic
Hamiltonian depends on an intrinsic second-level geometric object
rather than only on the closure of the first-level one. This is why
it is used in the present work as the primary diagnostic of the
breakdown of the standard sector. The important point is not that
$H_{22}$ must be the numerically largest post-standard quartic
contribution in every molecule, but that it is the first contribution
expected to break the linear representation transport that governs the
standard quartic and sextic sectors.

This role is specific. Purely vibrational channels such as
$H_{30}H_{30}$ do not introduce a new rotational tensorial level of
the same type, and mixed channels such as $H_{12}H_{30}$ still couple
anharmonicity to the first-level rotational object. For this reason
$H_{22}$ is the natural first candidate for the breakdown of linear
representation equivalence, even if other channels may later become
numerically larger in absolute value.

Whether $H_{22}$ becomes quantitatively important depends on the
molecule. The results section shows that it is small in systems such
as formaldehyde, where the standard perturbative picture is already
rather reliable, and larger in water, where the difference between
standard VPT-based and higher-level results is more pronounced. The
role of $H_{22}$ is therefore diagnostic rather than universal: it
signals the first failure of the standard linear sector, but does not
claim to exhaust the whole beyond-standard correction.

Beyond this point the tensor framework itself does not cease to apply.
What fails is the sufficiency of the minimal first-level linear sector.
Additional tensorial layers must then be introduced, and the simple
representation-equivalent transport picture of the standard sector no
longer closes by itself.

\subsection{First sextic analogue: linear response in \texorpdfstring{$\tau(H_{22})$}{tau(H22)}}

The same logic extends naturally to the sextic side. Within the
standard sextic geometry functional, the quartic tensor enters
explicitly. Therefore, once the standard quartic contribution is
replaced by
\begin{equation}
\tau = \tau_{\mathrm{std}} + \tau(H_{22}),
\end{equation}
the first post-standard sextic contribution is the term linear in
$\tau(H_{22})$.

Here $\tau(H_{22})$ denotes the quartic tensor representative
generated by the $H_{22}$ channel. The corresponding linear-response
term is the natural sextic analogue of the quartic channel $H_{22}$:
it is the first sextic contribution that probes the breakdown of the
standard linear sector while still remaining simple enough to be
treated as a diagnostic rather than as part of a full higher-order
theory.

This point should be read in strict analogy with the quartic case.
The quantity is not introduced as a complete post-standard sextic
correction, but as the first controlled way in which the new quartic
tensorial level feeds into the sextic problem. In that sense it
provides the sextic companion of the quartic $H_{22}$ diagnostic.

At the physical level, the reason cubic anharmonicity becomes the first
source of nonlinearity is that it couples vibrational modes in a way
that can no longer be reduced to the first-level geometric response of
the inverse inertia tensor. Once these cubic couplings enter
explicitly, new tensor combinations appear that do not close within the
minimal linear sector generated by $\mu_1$ and its standard geometric
descendants.

\subsection{Relation to earlier tensorial formulations}

The tensorial structure underlying centrifugal distortion is not new
in itself. It is already implicit in the classical formulations of
Watson and Mills and in later rotational-tensor treatments of
effective Hamiltonians.\cite{Watson1968,Watson1977,Mills1972,Kirschner1983}
What is new here is not the existence of
tensorial content behind Watson constants, but the use of that content
as the explicit operational layer for:
\begin{enumerate}
\item pseudoinverse reconstruction from spectroscopic parameter sets,
\item transport between reductions and principal-axis representations,
\item common treatment of standard quartic and standard sextic sectors,
\item identification of the first post-standard diagnostics.
\end{enumerate}

The present work therefore differs from earlier tensorial viewpoints
less by introducing a new kind of tensor than by making the tensor
representatives the center of the computational and interpretative
workflow. In this sense, the novelty lies in the combination of
tensor reconstruction, gauge fixing, representation transport, and
controlled identification of the first breakdown terms.

\subsection{Scope of the present work}

The present paper is therefore not a theory of the full VPT4 problem.
Its scope is deliberately narrower:
\begin{enumerate}
\item identify the standard linear quartic/sextic sector;
\item formulate its transport between representations and reductions;
\item isolate the first quartic and sextic diagnostics of its breakdown.
\end{enumerate}

The detailed tensor derivations, the explicit BCH/contact-transformation
coefficients, the quartic and sextic projection formulas, and the full
construction of the $H_{22}$ and $\tau(H_{22})$ diagnostics are
deferred to the appendices.
